\documentclass[11pt]{article}
\usepackage[a4paper,margin=1in]{geometry}
\usepackage{amsmath,amsthm,amssymb}
\usepackage{mathtools}
\usepackage{enumitem}
\usepackage{xcolor}

\newtheorem{theorem}{Theorem}[section]
\newtheorem{proposition}[theorem]{Proposition}
\newtheorem{lemma}[theorem]{Lemma}
\newtheorem{corollary}[theorem]{Corollary}
\theoremstyle{definition}
\newtheorem{definition}[theorem]{Definition}
\newtheorem{remark}[theorem]{Remark}

\newcommand{\R}{\mathbb{R}}
\newcommand{\Om}{\Omega}
\newcommand{\Fr}{\mathcal{A}}
\newcommand{\X}{\mathcal X}
\newcommand{\rstar}{\rho_{*}}
\newcommand{\rdel}{\rho_\delta}
\newcommand{\rhad}{\widehat\rho}
\newcommand{\Hn}{\mathcal H^{n-1}}
\newcommand{\dD}{\mathcal D}
\newcommand{\OPEN}{\textcolor{red}{\bf OPEN}}
\newcommand{\CLOSED}{\textcolor{teal}{\bf closed}}

\title{Route $\delta$ to the Sharp Quantitative Faber--Krahn Inequality:\\
       Step-by-Step Status, Remaining Obstructions, and Candidate Fixes}
\author{}
\date{\today}

\begin{document}
\maketitle

\begin{abstract}
This digest summarises the current state of Plan~2 / Route $\delta$,
which proves the sharp quantitative Faber--Krahn inequality
$\delta_{\rm FK}(\Om)\ge c(n)\Fr(\Om)^{2}$ via Saint--Venant stability
and a Kohler--Jobin transfer.  Each step of the proof chain is
condensed to its essential idea, its input, and its output, with
\CLOSED{} or \OPEN{} tags.  The remaining genuinely open question
collapses to a single integrated homothetic-trace estimate on
isoperimetrically bad levels.  We list six candidate routes for
closing it, with explicit critique of why each does or does not
appear to work.
\end{abstract}

\section{The target and the reduction}\label{sec:target}

\paragraph{Target (Faber--Krahn).}  For every open $\Om\subset\R^{n}$
of finite measure,
\[
   \Bigl(\tfrac{|\Om|}{|B_{1}|}\Bigr)^{2/n}
   \bigl(\lambda_{1}(\Om)-\lambda_{1}(B^{*})\bigr)
   \;\ge\; c_{\rm FK}(n)\,\Fr(\Om)^{2},
   \qquad\text{$B^{*}$ ball with }|B^{*}|=|\Om|.
\]
The exponent~$2$ is sharp; equality is approached by smooth
volume-preserving perturbations of the ball.

\paragraph{Reduction to Saint--Venant.}  Kohler--Jobin transfers
sharp Saint--Venant deficit stability to sharp Faber--Krahn:
if $\delta_{T}(\Om):=E(\Om)-E(B^{*})$ controls $\Fr(\Om)^{2}$
at rate~$\delta_{T}$, then $\delta_{\rm FK}$ also controls
$\Fr(\Om)^{2}$ at rate~$\delta_{\rm FK}$.  \CLOSED.  Reference:
\texttt{Final/KohlerJobinTransfer.tex} (Plan~1, imported).

\paragraph{Bounded reduction.}  The bounded sharp Saint--Venant
inequality on $\Om\subset B_{R}$ with $|\Om|=\omega_{n}$ implies the
global one via a Cicalese--Leonardi / BDV selection. \CLOSED.
Reference: \texttt{Final/BoundedReduction.tex}.

\paragraph{The bounded target.}
\[
   \Fr(\Om)^{2}\;\le\;C(n,R)\,\delta_{T}(\Om).
\]
This is what Route $\delta$ aims to prove.  Throughout the rest of
this digest, $|\Om|=\omega_{n}$, $\Om\subset B_{R}$, and $u=u_{\Om}$
denotes the variational torsion function.  Write
\[
   E_{\rho}:=\{u>t(\rho)\},\quad |E_{\rho}|=\omega_{n}\rho^{n},\quad
   F_{\rho}:=\rho^{-1}E_{\rho}\in\X,
\]
where $\X$ is the $L^{1}$-modulo-translations metric space.  Set
$w(\rho):=-t'(\rho)$, $d\mu:=w\,d\rho$, and
$\rdel:=1-\kappa\sqrt{\delta_{T}}$.

\section{The eight-step chain}\label{sec:chain}

\begin{center}
\renewcommand{\arraystretch}{1.2}
\begin{tabular}{|l|l|l|}
\hline
\textbf{Step} & \textbf{Block} & \textbf{Status} \\
\hline
1. Exact level-set deficit identity & \texttt{LevelSetIdentity} & \CLOSED \\
2. Metric first variation in $\X$ & \texttt{LimsupDeformation} (new) & \CLOSED \\
3. Same-centre Fusco--Julin package & \texttt{SameCentreFJ} (new) & \CLOSED \\
4. Trimmed velocity / good levels & \texttt{TrimmedVelocityRepair} & \CLOSED \\
4$'$. Unconditional velocity (all levels) & \texttt{UnconditionalVelocity} (new) & \CLOSED \\
5. Integrated action $D_{H}+D_{I}\Rightarrow d_{\X}^{2}+|\dot F|^{2}$ & \texttt{WeightedMetricTrace} & \OPEN $^{\dagger}$ \\
6. Uniform window trace & \texttt{WeightedMetricTrace} & \OPEN $^{\dagger}$ \\
6$'$. Good Borel endpoint $\rhad$ near $\rdel$ & \texttt{GoodEndpointTrace} (new) & \OPEN $^{\dagger}$ \\
7. Boundary-layer transfer from $E_{\rhad}$ to $\Om$ & \texttt{BoundaryLayerTransfer} & \CLOSED \\
8. Bounded $\Rightarrow$ global $\Rightarrow$ FK & \texttt{GlobalAssembly} & \CLOSED \\
\hline
\end{tabular}
\end{center}

$\dagger$~Steps 5--6 are open only modulo a \emph{single} integrated
bound on bad levels; see Section~\ref{sec:gap}.

\section{Step-by-step summary}\label{sec:step-by-step}

\subsection*{Step 1.\ Level-set deficit identity}\label{step:1}
\textbf{Statement.}
\[
   \frac{2\delta_{T}(\Om)}{|\Om|}
   =\frac{1}{|\Om|}\int_{0}^{\|u\|_{\infty}}
   \frac{D_{H}(t)+D_{I}(t)}{n^{2}\omega_{n}^{2/n}m(t)^{1-2/n}}\,dt,
\]
where $D_{H}(t)=\alpha(t)\beta(t)-P(E_{t})^{2}$ is the Cauchy--Schwarz
defect of $|\nabla u|$ on $\partial^{*}E_{t}$, and
$D_{I}(t)=P(E_{t})^{2}-P(B(E_{t}))^{2}$ is the squared
isoperimetric defect.\\
\textbf{Output.}  $\int_{\rstar}^{\rdel}(D_{H}+D_{I})\,d\mu\le C(n,\rstar)\delta_{T}$.\\
\textbf{Proof tools.}  BV coarea on reduced boundaries, Lipschitz
truncation for the flux identity
$\int_{\partial^{*}E_{t}}|\nabla u|\,d\Hn=|E_{t}|$, corrected
$I''=-1/\alpha$ for the rearranged primitive, P\'olya--Szeg\H{o}
endpoint convexity.

\subsection*{Step 2.\ Metric first variation (Deliverable~A)}\label{step:2}
\textbf{Statement.}  For a.e.\ $\rho$ and every bounded Borel centre
field $z_{\rho}$,
\[
   |\dot F_{\rho}|_{\X}
   \;\le\; \rho^{-n}\inf_{a\in\R^{n}}
   \int_{\partial^{*}E_{\rho}}|V_{\rho}-H_{z_{\rho},\rho}-a\cdot\nu_{\rho}|\,d\Hn,
\]
with $V_{\rho}=w/|\nabla u|$, $H_{z,\rho}(x)=(x-z)\cdot\nu_{\rho}(x)/\rho$.\\
\textbf{Proof idea.}  Direct coarea on the symmetric difference
$E_{\rho+h}\Delta T_{h}(E_{\rho})$, where $T_{h}$ is the affine
pull-back combining homothety $y\mapsto z+s^{-1}(y-z)$ and a chosen
translation $hb$.  No Reshetnyak or BV lower-semicontinuity.\\
\textbf{Status.}  \CLOSED.  Replaces the conditional upper-gradient
recovery hypothesis of the older \texttt{MetricFramework}.

\subsection*{Step 3.\ Same-centre Fusco--Julin (Deliverables~B+C)}\label{step:3}
\textbf{Statement.}  For every $E\subset K\subset\R^{n}$ of finite
perimeter with $|E|=\omega_{n}$ and $\dD(E)\le\eta_{0}(n,K)$, there
exists a single centre $z_{E}\in K$ with
\[
   \alpha_{FJ}(E,z_{E})^{2}+b_{FJ}(E,z_{E})^{2}\le C_{\rm scFJ}(n)\,\dD(E).
\]
Moreover the centre can be chosen \emph{Borel measurable} in $\rho$
along the torsion family.\\
\textbf{Proof idea.}  Joint functional
$\Phi_{E}(z):=\alpha_{FJ}^{2}+b_{FJ}^{2}$ on the compact admissible
$K$; LSC + compactness give attainment; the centres-close estimate
$|z_{F}-z_{FJ}|\le C\dD^{1/2}$ (via Fuglede / Cicalese--Leonardi
nearly-spherical) gives the bound; Castaing--Valadier provides Borel
selection.\\
\textbf{Status.}  \CLOSED.  Replaces the conditional theorems of
\texttt{FJCenterBridge}.

\subsection*{Step 4.\ Velocity defect}\label{step:4}
\textbf{Statement (unconditional, Theorem 2.2 of \texttt{UnconditionalVelocity}).}
For every a.e.\ finite-perimeter level,
\[
   \Bigl(\int_{\partial^{*}E_{\rho}}|V_{\rho}-\bar V_{\rho}|\,d\Hn\Bigr)^{\!2}
   \;\le\; \Bigl(\frac{P(B_{\rho})}{P(E_{\rho})}\Bigr)^{\!2} D_{H}(t(\rho))
   \;\le\; D_{H}(t(\rho)).
\]
\textbf{Proof idea.}  The weighted variance identity
$\int(V-\bar V)^{2}/V\,d\Hn=P^{*}D_{H}/P^{2}$ (closed form from
$w\alpha=P^{*}$ and $\beta=|E_{\rho}|$); Cauchy--Schwarz
$(\int|g|)^{2}\le(\int g^{2}/V)(\int V)$.\\
\textbf{Status.}  \CLOSED.

\subsection*{Step 5.\ Integrated action and kinetic energy}\label{step:5}
\textbf{Targets.}
\begin{align*}
   \text{(5a)}&\;\;\int_{\rstar}^{\rdel}\bigl(d_{\X}(F_{\rho},B_{1})^{2}+|\dot F_{\rho}|_{\X}^{2}\bigr)\,d\mu
      \;\le\; C\,\delta_{T},\\
   \text{(5b)}&\;\;\int_{\rstar}^{\rdel}|\dot F_{\rho}|_{\X}^{2}\,d\rho
      \;\le\; C\,\delta_{T}\quad\text{(Lebesgue form)}.
\end{align*}
\textbf{Status.}  On $G_{I}=\{D_{I}\le\eta_{I}\}$, the chain
[Step~2]+[Step~3]+[Step~4] gives the pointwise
$|\dot F_{\rho}|^{2}\le C(D_{H}+D_{I})$, hence \CLOSED{} on
$G_{I}$.  Integrated against $d\mu$, the level-set identity closes
the $G_{I}$ piece unconditionally.\\
The \emph{bad-isoperimetric} piece, $\int_{B_{I}}|\dot F|^{2}d\mu$,
is currently an explicit hypothesis ((eq:bad-I-action) of
\texttt{WeightedMetricTrace}).  \OPEN.

\subsection*{Step 6.\ Trace at the deficit endpoint}\label{step:6}
\textbf{Statement.}  $\sup_{\rho\in J^{*}}d_{\X}(F_{\rho},B_{1})^{2}\le C\delta_{T}$
on the order-one window $J^{*}=[\rdel-L^{*},\rdel]$.\\
\textbf{Proof idea.}  Markov on (5a) inside $J^{*}$ gives a single
$\rho_{0}\in J^{*}$ with $d^{2}\le K_{1}\delta_{T}$.  Kinetic
Cauchy--Schwarz with (5b) transports this to all of $J^{*}$.\\
\textbf{Status.}  \CLOSED{} conditional on (5b), which in turn
depends on Step~5's bad-level hypotheses.\\
\textbf{Variant (Deliverable~D).}  Replace ``trace at $\rdel$'' by
``trace at a Borel good level $\rhad\in[\rdel-C_{1}\delta_{T},\rdel]\cap
G_{I}\cap G_{\tau}$'' via the bad-set Lebesgue bounds
$\mathcal L^{1}(B_{\tau})\le C\delta_{T}$,
$\mathcal L^{1}(B_{I}\cap G_{\tau})\le C\delta_{T}$.  This cleans the
downstream interface but does \emph{not} eliminate the conditional
hypotheses of Step~5.

\subsection*{Step 7.\ Boundary-layer transfer}\label{step:7}
\textbf{Statement.}
$\Fr(\Om)\le\Fr(E_{\rhad})+\frac{2}{\omega_{n}}|\Om\setminus E_{\rhad}|$,
and $|\Om\setminus E_{\rhad}|\le\omega_{n}n(\kappa\sqrt{\delta_{T}}+C_{1}\delta_{T})
=O(\sqrt{\delta_{T}})$.\\
Squaring: $\Fr(\Om)^{2}\le C\delta_{T}+O(\delta_{T}^{3/2})$; sharp
exponent preserved. \CLOSED.

\subsection*{Step 8.\ Bounded $\Rightarrow$ FK}\label{step:8}
\textbf{Statement.}  Combine Step~7 with the bounded reduction (BDV /
Cicalese--Leonardi selection imported from Plan~1) and Kohler--Jobin
to obtain the global FK at the sharp exponent~2. \CLOSED.

\section{The single remaining obstruction}\label{sec:gap}

\textbf{The genuinely open question.}  The whole chain reduces to:
\emph{produce a bounded Borel centre field $\rho\mapsto z_{\rho}\in\R^{n}$ such that}
\begin{equation}\label{eq:hom-target}
   \int_{\rstar}^{\rdel}\Bigl(\int_{\partial^{*}E_{\rho}}|H_{z_{\rho},\rho}-\bar V_{\rho}|\,d\Hn\Bigr)^{\!2}d\mu(\rho)
   \;\le\;C(n,R,\rstar)\,\delta_{T}(\Om).
\end{equation}

\textbf{Why this is enough.}  Combining Steps~2, 4 (unconditional), and
\eqref{eq:hom-target} via the triangle inequality
$|V-H_{z}|\le|V-\bar V|+|H_{z}-\bar V|$ and squaring gives the full
integrated kinetic bound
$\int|\dot F_{\rho}|^{2}d\mu\le C\delta_{T}$ on \emph{all} levels.
In particular both bad-level kinetic hypotheses of \texttt{WeightedMetricTrace}
(\textsc{bad-I-action}, \textsc{bad-tau-action}) are discharged.

\textbf{Why this is hard.}  On the Fusco--Julin good set $G_{I}$, the
same-centre theorem (Step~3) yields
$\bigl(\int|H_{z_{\rho},\rho}-\bar V_{\rho}|\bigr)^{2}\le C\,\dD(E_{\rho})\le C'\mathcal I(\rho)$
pointwise, so \eqref{eq:hom-target} restricted to $G_{I}$ follows from
the level-set identity.  On the bad set
$B_{I}=\{\mathcal I>\eta_{I}\}$ the FJ inequality still holds
globally as $b_{FJ}(E,z_{FJ})^{2}\le C\dD(E)$, but \emph{the
Fraenkel and FJ centres no longer coincide}: a centre $z$ minimising
the joint $\alpha^{2}+b^{2}$ may not exist, and the radial-trace
identity
\[
   \int|H_{z}-\bar V|\le C[\,|E\Delta B(z)|+\int|\nu-e_{z}|^{2}\,]
\]
loses on one side or the other at any single $z$.

\textbf{Quantitative obstruction.}  Pointwise the perimeter
$P(E_{\rho})$ can be unboundedly larger than $P(B_{\rho})$ on bad
levels (a torsion superlevel set may be very rough), and the
homothetic integrand $|H_{z}-\bar V|$ has uncontrolled $L^{2}$ norm
on $\partial^{*}E_{\rho}$.  The $L^{1}$ norm is controlled by the
asymmetry and the normal oscillation but not in a form that integrates
to $C\delta_{T}$ after squaring.

\section{Candidate routes for closing \eqref{eq:hom-target}}\label{sec:fixes}

\subsection{Route I: centroid $+$ space--time $\Pi$-identity}
\label{ssec:routeI}

\textbf{Idea.}  Set $z_{\rho}:=\bar z_{\rho}=$ centroid of $E_{\rho}$.
The divergence identity
\[
   \int_{\partial^{*}E_{\rho}}|\nu_{E_{\rho}}-e_{\bar z_{\rho}}|^{2}d\Hn
   =2\bigl(P(E_{\rho})-P(B_{\rho})\bigr)+2\,\Pi(E_{\rho},\bar z_{\rho})
\]
of \texttt{SpaceTimeIdentity} converts the normal oscillation into
a quadratic form $\Pi$, and the pointwise slice-rearrangement
inequality $|E_{\rho}\Delta B_{\rho}(\bar z_{\rho})|^{2}\le C\,\Pi$
controls the Fraenkel term by $\Pi$ too.\\
\textbf{Why it might work.}  \emph{If} the conditional integrated
estimate $\int\Pi\,d\mu\le C\delta_{T}$ (``$\Pi$-int'') can be proved,
then $\int\int|\nu-e_{\bar z}|^{2}d\Hn d\mu\le C\delta_{T}$ globally
(good and bad levels both), and Cauchy--Schwarz then bounds
\eqref{eq:hom-target} at the centroid centre.\\
\textbf{Why it might not.}  $\Pi$-int is currently an unproved
structural conjecture (the abstract Talenti shortcut was invalidated
in the May 14 audit).  More seriously, even granted $\Pi$-int, the
\emph{squared} integral $\int(\int|\nu-e|^{2}d\Hn)^{2}d\mu$ is not
bounded by $\delta_{T}$ via naive Cauchy--Schwarz: one only obtains
$\sqrt{\delta_{T}}$, since
$(\int|\nu-e|^{2})^{2}\le 4P\int|\nu-e|^{2}$ and
$\int P(E_{\rho})^{2}d\mu\le\omega_{n}^{2}$ is constant.  A second
deficit-aware inequality on $P$ is needed.

\textbf{Critique.}  The bottleneck is not really $\Pi$ but the
$L^{2}\to L^{2}$ stability of $\Hn|\partial^{*}E_{\rho}$ across $\mu$;
$\Pi$ controls the $L^{1}$ mass, and the squared form is genuinely
new information.  Route~I would close the bound \emph{up to a
$\sqrt{\delta_{T}}$ exponent loss}, not at the sharp~$\delta_{T}$ rate.

\subsection{Route II: Fraenkel centre $+$ direct radial trace}
\label{ssec:routeII}

\textbf{Idea.}  Use the Fraenkel centre $z_{F}(E_{\rho})$, at which
$|E_{\rho}\Delta B_{\rho}(z_{F})|^{2}\le C_{F}\mathcal I(\rho)$
holds \emph{globally} (Figalli--Maggi--Pratelli, no smallness).
Combine with the radial trace from \texttt{FJCenterBridge}.\\
\textbf{Why it might work.}  Avoids any FJ centre, so works on bad
levels.  The Fraenkel asymmetry is the only quantity needed for the
$\alpha$-piece.\\
\textbf{Why it might not.}  The normal oscillation $\int|\nu-e_{z_{F}}|^{2}$
at the \emph{Fraenkel} centre is uncontrolled on bad levels: FJ
controls the oscillation at $z_{FJ}$, not at $z_{F}$, and the
``centres-close'' bound $|z_{F}-z_{FJ}|\le C\dD^{1/2}$ holds
\emph{only on good levels} (via nearly-spherical).  At a bad level
with comparable Fraenkel and FJ values but distinct centres, the
oscillation at $z_{F}$ can be of order $P(E_{\rho})$, which is the
$L^{\infty}$-trivial bound.

\textbf{Critique.}  This route would require a non-trivial new
``centres are close'' theorem in the absence of nearly-spherical
structure; this seems not to be in the literature and probably
demands its own large argument.

\subsection{Route III: Schwarz/Steiner symmetrisation}\label{ssec:routeIII}

\textbf{Idea.}  Replace $E_{\rho}$ by its Schwarz symmetrisation
$E_{\rho}^{\#}$ (a ball).  Then by P\'olya--Szeg\H{o}, the torsion
energy decreases under symmetrisation; the deficit
$\delta_{T}(\Om^{\#})=0$, so all level-set defects vanish for
$\Om^{\#}$.  The difference between $E_{\rho}$ and $E_{\rho}^{\#}$ is
controlled by the Saint--Venant deficit.\\
\textbf{Why it might work.}  Symmetrisation is the standard tool for
\emph{relating} $E_{\rho}$ to a ball; in principle the deficit
$\delta_{T}$ itself is the price of asymmetry, so a quantitative
P\'olya--Szeg\H{o} would give a level-by-level bound.\\
\textbf{Why it might not.}  Quantitative P\'olya--Szeg\H{o} stability
(Cianchi--Esposito--Fusco--Trombetti~2011, Brock~2000) gives a bound
like $\|\nabla u\|_{2}^{2}-\|\nabla u^{\#}\|_{2}^{2}\ge c\cdot$~asymmetry,
which involves $\|\nabla u\|$, not the level-set defect.  To
translate into the $\X$-metric on individual levels one would still
need a slice-by-slice argument, and at the slice level we are back
to the same homothetic-trace question.

\textbf{Critique.}  Symmetrisation reproves Saint--Venant deficit
bounds but does not produce a centre field with control on bad
levels.  Useful as a sanity check but not, by itself, a closure.

\subsection{Route IV: Cicalese--Leonardi / BDV selection at the level}
\label{ssec:routeIV}

\textbf{Idea.}  On a bad level $E_{\rho}$, replace it by a ``selection
representative'' $\widetilde E_{\rho}$ with small deficit, then run
FJ on $\widetilde E$ and transfer back to $E_{\rho}$ with a bounded
error.\\
\textbf{Why it might work.}  This is exactly the philosophy of the
Cicalese--Leonardi quantitative isoperimetric theorem, which gave the
sharp FMP exponent~2 from a selection step.  Adapting it to torsion
levels would automatically handle bad levels at sharp rate.\\
\textbf{Why it might not.}  The Cicalese--Leonardi selection is a
\emph{compactness} argument: for sequences with $\dD\to 0$, after
translation one extracts a convergent subsequence to the ball.  It
does \emph{not} provide a quantitative pointwise replacement
$E_{\rho}\rightsquigarrow\widetilde E_{\rho}$ at a single level
without small deficit.  Adapting the argument to bounded but not
small deficit would amount to extending the regularity theory of
quasi-minimal surfaces, which is a major undertaking beyond Route
$\delta$'s scope.

\textbf{Critique.}  Route~IV is the cleanest in principle but
hardest in practice: it would require new selection theorems for
torsion superlevel sets, not direct application of existing FMP/FJ
machinery.

\subsection{Route V: optimal transport $W_{1}$-control}\label{ssec:routeV}

\textbf{Idea.}  Replace $|E\Delta B(z)|$ and $\int|\nu-e_{z}|^{2}$ by
$W_{1}(\mathbf 1_{E},\mathbf 1_{B(z)})$, the Kantorovich--Wasserstein
distance.  For convex sets there are quantitative Bonnesen-type
inequalities relating $W_{1}$ to the isoperimetric deficit.\\
\textbf{Why it might work.}  Optimal transport gives a single
unified quantity controlling both ``shape'' and ``boundary
orientation'' simultaneously; in principle $W_{1}\le C\sqrt{\dD}$
covers both Fraenkel and oscillation in one inequality.\\
\textbf{Why it might not.}  No quantitative
$W_{1}(\mathbf 1_{E},\mathbf 1_{B})\le C\dD^{1/2}$ inequality is
known for non-convex sets at the sharp rate.  Existing Wasserstein
isoperimetric stability (Figalli--Indrei) gives weaker rates, and the
relationship $W_{1}\Leftrightarrow$~radial-trace integrand is not
direct: the radial trace measures normal alignment, whereas $W_{1}$
measures bulk transport.

\textbf{Critique.}  Aesthetically attractive but the dictionary
between $W_{1}$ and the radial trace is not established at the
required sharp rate.

\subsection{Route VI: pointwise everywhere via algebraic identity}
\label{ssec:routeVI}

\textbf{Idea.}  Look for a single closed-form identity of the type
$\int(H_{z}-\bar V)^{2}/V\,d\Hn\;\le\;C(n,R)\,(D_{H}+D_{I})$
holding pointwise on every level, parallel to the velocity identity
$\int(V-\bar V)^{2}/V=P^{*}D_{H}/P^{2}$.  This would give the
homothetic analogue of Theorem~2.2 of \texttt{UnconditionalVelocity}.\\
\textbf{Why it might work.}  The velocity identity is genuinely
algebraic, requiring only the flux $w\alpha=P^{*}$ and
$\beta=|E_{\rho}|$.  A parallel identity for $H_{z}$ would use the
divergence identity $\int(x-z)\cdot\nu=n|E_{\rho}|$ and the Reilly
or Bochner formulas for torsion.  Plausible in principle.\\
\textbf{Why it might not.}  Concrete attempts produce closed forms
involving $\int|x-z|^{2}d\Hn$ or $\int((x-z)\cdot\nu)^{2}$, which
have no clean divergence representation (the corresponding vector
field is not exact).  $D_{I}$ itself does not appear naturally in
these computations; $D_{I}$ is intrinsically a global
``perimeter-squared minus ball-perimeter-squared'' quantity, while
the homothetic integrand is local.  No identity is currently
visible.

\textbf{Critique.}  This is the route most analogous to the proved
velocity bound (Step~4$'$).  The reason it does not directly
generalise is that $V_{\rho}$ comes from the gradient (a vector
field with a specific divergence), while $H_{z}$ comes from the
\emph{position}~$x$ (also a vector field, but its second moments
are not local on the boundary).  This is a real geometric obstruction.

\section{Worked attempts that fail — and what they show}\label{sec:failed-attempts}

We record concretely what happens when each natural Cauchy--Schwarz is
applied to the homothetic integrand, so the geometric obstruction is
visible.

\subsection{Attempt~A: weighted Cauchy--Schwarz in $V_\rho$}

Mirror the proof of Theorem~4.1 of \texttt{UnconditionalVelocity}.
\[
   \Bigl(\int|H_{z}-\bar V|\,d\Hn\Bigr)^{2}
   \;\le\;\Bigl(\int\frac{(H_{z}-\bar V)^{2}}{V_{\rho}}\,d\Hn\Bigr)\Bigl(\int V_{\rho}\,d\Hn\Bigr).
\]
The second factor is $P(B_\rho)=P^{*}$. The first factor expands to
\[
   \int(H_{z}-\bar V)^{2}\frac{|\nabla u|}{w}\,d\Hn
   =\frac{1}{w}\int(H_{z}-\bar V)^{2}\,|\nabla u|\,d\Hn,
\]
a $|\nabla u|$-weighted $L^{2}$-norm of $H_{z}-\bar V$.  This is
finite (since $H_{z}$ is bounded on $\partial^{*}E_{\rho}\subset B_{R}$)
but not bounded by any deficit quantity: a level surface with constant
$|\nabla u|=\bar f$ has $D_{H}=0$ but generic $H_{z}-\bar V\ne 0$, so
the LHS is not even nonzero where $D_{H}$ is.

\emph{What this shows.}  $V_{\rho}$-weighting is too gradient-specific.
The homothetic integrand is geometric, not gradient-geometric.

\subsection{Attempt~B: unweighted Cauchy--Schwarz}

\[
   \Bigl(\int|H_{z}-\bar V|\,d\Hn\Bigr)^{2}
   \;\le\; P(E_{\rho})\,\int(H_{z}-\bar V)^{2}\,d\Hn.
\]
For $z=\bar z_{\rho}$ the centroid, expand
$(H_{z}-\bar V)^{2}=H_{z}^{2}-2\bar V H_{z}+\bar V^{2}$ and integrate:
$\int H_{z}=P^{*}$, $\bar V P=P^{*}$, so
\[
   \int(H_{z}-\bar V)^{2}\,d\Hn
   =\int H_{z}^{2}\,d\Hn-\frac{P^{*\,2}}{P}.
\]
The integral $\int H_{z}^{2}=\rho^{-2}\int((x-z)\cdot\nu)^{2}\,d\Hn$
has no clean divergence representation: the candidate vector field
$X(x):=((x-z)\cdot\nu)(x-z)$ requires the normal $\nu$ off
$\partial^{*}E_{\rho}$.  Bound:
$((x-z)\cdot\nu)^{2}\le|x-z|^{2}$, so
$\int H_{z}^{2}\le ((R+|z|)/\rstar)^{2}\,P(E_{\rho})/\rho^{2}$.  Hence
\[
   \Bigl(\int|H_{z}-\bar V|\Bigr)^{2}
   \;\le\;C(R,\rstar)\,P(E_{\rho})^{2}.
\]
The square of the perimeter is unbounded on bad levels; integrated,
$\int P(E_{\rho})^{2}d\mu\le\omega_{n}^{2}$ is a \emph{constant}, not a
$\delta_{T}$-vanishing quantity.

\emph{What this shows.}  Without an additional deficit-aware bound on
$\int H_{z}^{2}$, naive Cauchy--Schwarz only reaches a constant, not
$O(\delta_{T})$.

\subsection{Attempt~C: split into radial and angular contributions}

Decompose $H_{z}-1=(|x-z|/\rho)(\nu\cdot e_{z}-1)+(|x-z|/\rho-1)$.  The
first piece has magnitude $\le (R/\rstar)|\nu-e_{z}|^{2}/2$, the
second has magnitude $w_{z}:=||x-z|/\rho-1|\le W(R,\rstar)$.  Then
\[
   \int|H_{z}-1|\,d\Hn
   \;\le\;\frac{R+|z|}{2\rho}\int|\nu-e_{z}|^{2}\,d\Hn
   \;+\;\int w_{z}\,d\Hn.
\]
The first piece, at $z=z_{FJ}$, is $\le C\mathcal I/P^{*}$ globally
(unconditional Fusco--Julin) — but, squared and integrated, only
yields $\int(\mathcal I/P^{*})^{2}d\mu\le\eta_{I}\int\mathcal I\,d\mu/P^{*\,2}\le C\delta_{T}$
\emph{on the good set}, and is unbounded on the bad set since
$\mathcal I$ is not pointwise bounded above.

The second piece is bounded by the radial trace
$\int w_{z}\le(n/\rstar)|E\Delta B_{\rho}(z)|+(W/2)\int|\nu-e_{z}|^{2}$,
so its squared integrated form inherits the same obstruction.

\emph{What this shows.}  Splitting helps separate the two
contributions but reduces to the same difficulty: on bad levels,
$\mathcal I$-pointwise control is missing.

\subsection{Attempt~D: choose $z$ optimal at each level}

Define $z_{\rho}^{*}:=\arg\min_{z\in K}\bigl(\int|H_{z,\rho}-\bar V_{\rho}|\,d\Hn\bigr)^{2}$.
By Castaing--Valadier this is a Borel selection (the functional is
continuous in $z$ and Borel in $\rho$).  Then
\[
   \int_{\rstar}^{\rdel}\Bigl(\int|H_{z_{\rho}^{*},\rho}-\bar V_{\rho}|\Bigr)^{2}d\mu(\rho)
   \;=\;\int_{\rstar}^{\rdel}\inf_{z}\Bigl(\int|H_{z}-\bar V|\Bigr)^{2}d\mu(\rho).
\]
\emph{What this shows.}  The Borel-selection issue is not the
obstruction; the obstruction is the pointwise question
\[
   \inf_{z\in K}\Bigl(\int|H_{z,\rho}-\bar V_{\rho}|\,d\Hn\Bigr)^{2}
   \;\le\;?
\]
for which no proved upper bound by $D_{H}+D_{I}$ is known.  This is
exactly $(\star)$ in the next section.

\section{What kind of input \emph{would} close \eqref{eq:hom-target}?}
\label{sec:wanted}

The cleanest sufficient input is a global pointwise inequality on
every a.e.\ finite-perimeter level:
\[
   (\star)\qquad\inf_{z\in K}\Bigl(\int|H_{z,\rho}-\bar V_{\rho}|\,d\Hn\Bigr)^{\!2}
   \;\le\;C(n,R,\rstar)\bigl(D_{H}(t(\rho))+D_{I}(t(\rho))\bigr).
\]
By the level-set identity, $(\star)$ implies \eqref{eq:hom-target}
and closes the entire route at the sharp exponent.

Equivalently, a weaker but still sufficient input is the
integrated form:
\[
   (\star\star)\qquad
   \int_{\rstar}^{\rdel}\inf_{z}\Bigl(\int|H_{z,\rho}-\bar V_{\rho}|\,d\Hn\Bigr)^{\!2}d\mu(\rho)
   \;\le\;C\,\delta_{T}(\Om).
\]

Routes I--VI each address $(\star\star)$ with a different choice of
$z_{\rho}$ and a different identification of the integrand.  The
crucial structural fact is that the optimum $z_{\rho}$ in
$(\star\star)$ should depend on the geometry of $E_{\rho}$, and
the level-set identity should be \emph{evaluated} at this $z_{\rho}$
in a way that mirrors how $D_{H}$ already controls the velocity
defect.

\section{Other consequences if the gap is closed}\label{sec:if-closed}

Closing \eqref{eq:hom-target} would:
\begin{enumerate}[label=(\arabic*)]
\item Discharge the conditional Hypothesis~3.6 of \texttt{WeightedMetricTrace};
\item Make the entire chain Steps~1--8 \emph{unconditional} except
for the imported BDV reduction and Kohler--Jobin (both already
unconditional in Plan~1);
\item Yield the sharp Faber--Krahn inequality
$\delta_{\rm FK}(\Om)\ge c(n)\Fr(\Om)^{2}$ for every open
$\Om\subset\R^{n}$ of finite measure, with no smoothness, convexity,
or smallness assumption beyond those already in the imports;
\item Provide an explicit polynomial bound on $c(n)$ in dimension
(via the polynomial constants tracked in
\texttt{GlobalAssembly}).
\end{enumerate}

\section{Concrete suggestions for next attempts}\label{sec:suggestions}

\begin{enumerate}[label=(S\arabic*)]
\item Push Route~I by proving the integrated squared-oscillation
bound
$\int(\int|\nu-e_{\bar z}|^{2})^{2}d\mu\le C\delta_{T}$ directly,
exploiting that on bad levels both factors of $\int|\nu-e_{\bar z}|^{2}$
already carry $\Pi$.
\item Push Route~II by proving a quantitative ``centres-close'' on bad
levels: under what \emph{geometric} hypothesis is
$|z_{F}-z_{FJ}|\le C\dD^{1/2}$ valid outside the nearly-spherical
regime?
\item Push Route~IV by adapting the Cicalese--Leonardi selection to
\emph{a single level} at finite (not small) deficit; this would
require a localised version of the selection principle.
\item Investigate Route~VI by writing the divergence identity for the
candidate vector field
$X(x):=(H_{z}(x)-\bar V_{\rho})((x-z)/|x-z|)$ on
$\partial^{*}E_{\rho}$, and computing the bulk integrand pointwise.
\end{enumerate}

\section{Dependency map of the route}\label{sec:dep-map}

\[
\begin{array}{rclcl}
\text{Step 1: \texttt{LevelSetIdentity}}
   & \longrightarrow & \int(D_{H}+D_{I})d\mu\le C\delta_{T} & & \CLOSED \\
\text{Step 2: \texttt{LimsupDeformation}}
   & \longrightarrow & |\dot F_{\rho}|\le \rho^{-n}\inf_{a}\int|V-H_{z}-a\nu|d\Hn & & \CLOSED \\
\text{Step 3: \texttt{SameCentreFJ}}
   & \longrightarrow & \alpha_{FJ}(E,z_{E})^{2}+b_{FJ}(E,z_{E})^{2}\le C\dD(E) & & \CLOSED \\
\text{Step 4: \texttt{TrimmedVelocityRepair}}
   & \longrightarrow & \text{good-level }(\int|V-\bar V|)^{2}\le C D_{H} & & \CLOSED \\
\text{Step 4$'$: \texttt{UnconditionalVelocity}}
   & \longrightarrow & \text{all-level }(\int|V-\bar V|)^{2}\le D_{H} & & \CLOSED \\
\fbox{\eqref{eq:hom-target}: integrated homothetic-trace}
   & \longrightarrow & \int(\int|H_{z}-\bar V|)^{2}d\mu\le C\delta_{T} & & \OPEN \\
\text{Step 5: \texttt{WeightedMetricTrace}}
   & \Longleftarrow & \text{Steps 1, 2, 4$'$, and \eqref{eq:hom-target}} & & \\
\text{Step 6: trace at endpoint}
   & \Longleftarrow & \text{Step 5} & & \\
\text{Step 7: \texttt{BoundaryLayerTransfer}}
   & \longrightarrow & \Fr(\Om)\le \Fr(E_{\rhad})+(2/\omega_{n})|\Om\setminus E_{\rhad}| & & \CLOSED \\
\text{Step 8: \texttt{GlobalAssembly}}
   & \longrightarrow & \text{Bounded SV}\Rightarrow\text{FK} & & \CLOSED \\
\end{array}
\]

\medskip
The single \OPEN{} input \eqref{eq:hom-target} propagates downward to
Steps~5 and~6.  Everything else in the chain is unconditional in the
current state of the proof.  All inputs to \eqref{eq:hom-target} are
local-to-a-single-level and depend only on the geometry of
$E_{\rho}\subset B_{R}$ (a finite-perimeter set), not on the deficit
$\delta_{T}(\Om)$.

\section{Cost analysis: what each candidate route requires}\label{sec:cost}

\begin{center}
\renewcommand{\arraystretch}{1.3}
\begin{tabular}{|l|p{5.5cm}|c|c|}
\hline
\textbf{Route} & \textbf{New input needed} & \textbf{Rate at goal} & \textbf{Difficulty} \\
\hline
I (centroid + $\Pi$)
   & Prove $\int\Pi d\mu\le C\delta_{T}$ AND $\int(\int|\nu-e|^{2})^{2}d\mu\le C\delta_{T}$
   & sharp $\delta_{T}$ if both
   & medium-hard \\
II (Fraenkel centre)
   & Prove $|z_{F}-z_{FJ}|\le C\dD^{1/2}$ outside the nearly-spherical regime
   & sharp $\delta_{T}$ if input proved
   & hard \\
III (symmetrisation)
   & Quantitative slice-wise P\'olya--Szeg\H{o}
   & weaker than $\delta_{T}$
   & medium \\
IV (selection)
   & Adapt Cicalese--Leonardi to a single bounded-deficit level
   & sharp $\delta_{T}$
   & very hard \\
V (optimal transport)
   & Sharp $W_{1}\le C\dD^{1/2}$ for non-convex sets
   & sharp $\delta_{T}$ if input proved
   & very hard \\
VI (algebraic identity)
   & A closed form
     $\int(H_{z}-\bar V)^{2}/V d\Hn\le C(D_{H}+D_{I})$
   & sharp $\delta_{T}$
   & hard \\
\hline
\end{tabular}
\end{center}

The "Rate at goal" column records the best achievable exponent of
$\delta_{T}$ in $\Fr(\Om)^{2}\le ?$ at the end of the route if the
listed input is granted.  Sharp~$\delta_{T}$ corresponds to the
sharp exponent~$2$ for $\Fr$.

\section{Why this gap, and why now}\label{sec:why}

The proof chain prior to the current pass had \emph{four} conditional
inputs:
\begin{enumerate}[label=(\arabic*)]
\item the upper-gradient recovery hypothesis (closed in Deliverable~A);
\item the same-centre Fusco--Julin theorem (closed in B);
\item the Borel measurable centre selection (closed in C);
\item the bad-level kinetic action bounds.
\end{enumerate}
Three of these are now structural theorems with proofs (algebraic,
quantitative-FJ, and measurable-selection respectively).  The fourth
is a \emph{quantitative geometric} statement about how rough a
torsion super-level set can be when its perimeter is large.

The reason the fourth is more obdurate is that the first three are
\emph{universal} (they apply to general finite-perimeter sets and
torsion functions, without reference to deficit), whereas the fourth
asks for a $\delta_{T}$-vanishing bound on a quantity that, at any
single bad level, can be arbitrarily large.  The required smallness
must come from a \emph{global} balance: bad levels are rare in
$\mu$-measure, and a refined Cauchy--Schwarz must exploit this
rareness.  Existing inequalities are not yet sharp enough to do so
at the $\delta_{T}$ rate after squaring.

\section{Summary}\label{sec:summary}

The Plan 2 / Route $\delta$ proof is closed except for a single
integrated homothetic-trace estimate, \eqref{eq:hom-target}.  Six
candidate routes are visible; none currently closes it at the sharp
$\delta_{T}$ rate.  Routes I and~VI are the most parallel to the
existing closed steps (algebraic identities + level-set deficit);
Route IV is the most general but requires new selection theory.  The
mathematical question is sharp and decoupled from Fusco--Julin; it
is the unique residual gap in an otherwise complete chain.

The way forward is to combine: (a) a deficit-aware refinement of the
Cauchy--Schwarz step in Route~I, (b) a quantitative ``centres-close''
theorem outside the nearly-spherical regime for Route~II, or (c) a
closed-form algebraic identity for the homothetic integrand
analogous to the velocity identity for Route~VI.  Each represents a
self-contained mathematical problem that does not depend on the rest
of the chain.

\end{document}
